\chapter{\prjname{FPGAlix}}
\label{chap:fpgalix-intro}

This part covers the realization of \prjname{FPGAlix}, the Linux-embedded variant of the video acquisition and processing system introduced in the \nameref{chap:preface}. The chapters that follow proceed in the order the work itself was carried out: setting up the development environment on the host \term{PC}, building the hardware side of the system, and assembling around it a working embedded system based on a \term{Linux} kernel and a \gterm{buildroot}{Buildroot} environment.\footnote{\glsentrydesc{buildroot}} Once the system boots and the network connection is verified, the \term{DE1-SoC}'s 7-segment displays are used to show the board's \term{IP} address. Finally, the part closes with a driver for the camera over the custom hardware path and an application that processes the video stream and streams the result over the network via \gterm{rtsp}{RTSP}.\footnote{\glsentrydesc{rtsp}}

\section{Quick Start}
\label{sec:fpgalix-quick-start}

Readers who want the finished system running immediately, without going through every step described in this part, can start here instead. A prebuilt SD card image is available: it already contains the bootloader, kernel, root filesystem and \gterm{fpga}{FPGA}\footnote{\glsentrydesc{fpga}} bitstream, letting the \term{DE1-SoC} boot straight into a running system, without building anything from source.

\subsection{Hardware needed}

\begin{itemize}
    \item A \term{Terasic DE1-SoC} developer board.
    \item An \gterm{ov7670}{OV7670}\footnote{\glsentrydesc{ov7670}} camera module.
    \item The custom camera adapter (transmitter and receiver boards, connected by a ribbon cable), wiring the camera to the \term{DE1-SoC}'s \term{GPIO\_0} header.
    \item A microSD card, at least 8 GB.
    \item An \term{Ethernet} cable, for network access.
    \item A \term{PC} on the same network, to reach the board over \gterm{ssh}{SSH}\footnote{\glsentrydesc{ssh}} and view the video stream.
    \item A micro-\term{USB} cable, for the board's serial console. Optional: only needed to log in as root directly, or for troubleshooting.
    \item A power supply for the board.
\end{itemize}

\subsection{Preparing the SD card}

\begin{enumerate}
    \item Download the SD card image\footnote{See Chapter~\ref{chap:fpgalix-downloads} for where to obtain it.} and extract it.
    \item Write the extracted image to the microSD card (at least 8 GB) with any disk-imaging tool, e.g. \cmdpath{dd} on \term{Linux} or \gterm{balenaetcher}{balenaEtcher}\footnote{\glsentrydesc{balenaetcher}} on Windows/Mac.
    \item Insert the card into the \term{DE1-SoC}.
\end{enumerate}

\subsection{Board configuration}

Before powering on, set the \term{DE1-SoC}'s \term{SW10} DIP switch, on the underside of the board, so that switches 1 through 5 are ON (\term{SW10} uses negative logic: ON means logic 0). This lets the \gterm{hps}{HPS}\footnote{\glsentrydesc{hps}} configure the \gterm{fpga}{FPGA} fabric from the SD card at boot.\footnote{See Chapter~\ref{chap:fpgalix-board-bring-up}.}

\subsection{Booting and logging in}

\begin{warningbox}
\textbf{Hardware bug:} do not leave the camera in \term{RESET} or \term{POWER-DOWN} for extended periods.\footnotemark
\end{warningbox}
\footnotetext{See Part~\ref{part:appendices}, Section~\ref{sec:camera-adapter-reset-powerdown}.}

Connect the \gterm{ov7670}{OV7670} camera to the \term{DE1-SoC} board's \term{GPIO\_0} header via the camera adapter,\footnote{See Part~\ref{part:appendices}, Chapter~\ref{app:camera-adapter-board}.} connect the \term{Ethernet} cable, insert the SD card, and power on the board.

The board's \term{IP} address is shown automatically, a few seconds later, on the 7-segment displays, cycled one octet at a time by \cmd{hexip_daemon}.\footnote{See Chapter~\ref{chap:fpgalix-hexip-daemon}.}

Two ways to log in, default password |1q2w3e4r| for both accounts:

\begin{itemize}
    \item over \gterm{ssh}{SSH}, once the IP is known, as \sysid{wheel}, then \cmdpath{su} for anything that needs root (loading \cmd{camera_driver_2}, accessing hardware registers directly):

\begin{shellcode}
ssh wheel@<board-ip>
\end{shellcode}

    \item over the serial console (115200 baud, 8N1, e.g. \cmdpath{picocom}), directly as \sysid{root}; no network needed for this path.
\end{itemize}

\subsection{Running the software}

\cmd{hexip_daemon} is already running once the board boots; no action needed. The prebuilt SD card image also already contains \cmd{camera_driver_2} and \cmd{vision_core}, built and placed in \sysid{wheel}'s own home directory (\extpath{/home/wheel/camera_driver_2/}, \extpath{/home/wheel/vision_core/}); no building or transferring files over \cmdpath{scp} is needed here. Building from source and deploying over \cmdpath{scp}, covered in full in later chapters, only matters after changing their source.

Log in as \sysid{wheel} (or \sysid{root} from the serial console, cd-ing to \extpath{/home/wheel} first), then bring the camera up and load the driver:

\begin{shellcode}
cd camera_driver_2
su                              # Not needed if already logged in as root.
./resetCtrl.sh                  # option 2, deassert reset
./ov7670_ctrlif_set.sh          # select Bayer mode (1)
./ov7670_dataif_ctrl.sh         # option 1, enable frame acquisition
insmod fpgalix_camera.ko        # load camera driver module
\end{shellcode}

Then run \cmd{vision_core}, still as \sysid{root}, from the sibling directory:

\begin{shellcode}
cd ../vision_core
./vision_core
\end{shellcode}

It prompts interactively; every prompt accepts its default by pressing Enter, except the source, which must be set explicitly to \hw{ov7670}:

\begin{shellcode}
Source [webcam/ov7670, default: webcam]: ov7670
Encoding [mjpeg/h264/raw, default: mjpeg]:
\end{shellcode}

\hw{ov7670} is mandatory on the board, the \ccode{webcam} source being only for the native \term{Intel} (\cmdpath{make intel}) build; \ccode{MJPEG}, left at its default by pressing Enter, is what this project recommends for full-color output here, since the \gterm{hps}{HPS} has no hardware video encoder.

Once running, view the stream on a PC:

\begin{shellcode}
ffplay -fflags nobuffer -flags low_delay -framedrop rtsp://<board-ip>:8554/stream
\end{shellcode}

\section{Repository layout}

The \prjname{FPGAlix} repository \cite{FPGALIX-REPO}, which is hosted on \gterm{github}{GitHub}\footnote{\glsentrydesc{github}}, is laid out as follows:

\begin{itemize}
    \item \repopath{hw/}: \gterm{fpga}{FPGA} hardware project (\gterm{quartusprime}{Quartus Prime}\footnote{\glsentrydesc{quartusprime}} + \gterm{platformdesigner}{Platform Designer}\footnote{\glsentrydesc{platformdesigner}}).
    \begin{itemize}
        \item \repopath{camera_adapter/}: custom \gterm{pcb}{PCB}\footnote{\glsentrydesc{pcb}} bridging the \gterm{ov7670}{OV7670} to the \term{DE1-SoC}'s \term{GPIO\_0} header; the board layouts themselves are mirrored in the shared documentation folder (Chapter~\ref{chap:fpgalix-downloads}).
        \item \repopath{my_vhdl/}: custom \gterm{platformdesigner}{Platform Designer} \gterm{vhdl}{VHDL}\footnote{\glsentrydesc{vhdl}} components written for this project.
        \item \repopath{quartus/}: \term{Quartus} project, \repopath{soc_system.qsys}, top-level \gterm{vhdl}{VHDL}, pin assignments.
    \end{itemize}
    \item \repopath{sw/}: application software running on the board's embedded \term{Linux} (\term{C/C++}), one folder per component.
    \begin{itemize}
        \item \repopath{camera_driver_2/}: \gterm{v4l2}{V4L2}\footnote{\glsentrydesc{v4l2}} kernel driver for the \gterm{ov7670}{OV7670}.
        \item \repopath{hexip_daemon/}: displays the board's IP on the 7-segment displays.
        \item \repopath{vision_core/}: captures, filters and streams the camera feed over \gterm{rtsp}{RTSP}.
    \end{itemize}
    \item \repopath{repos/}: git submodules (\gterm{uboot}{U-Boot}\footnote{\glsentrydesc{uboot}}, \term{Linux} kernel, \gterm{buildroot}{Buildroot}, \gterm{ghrd}{GHRD}\footnote{\glsentrydesc{ghrd}}, \repopath{sopc2dts}, \gterm{altera}{Altera}\footnote{\glsentrydesc{altera}} \gterm{fpga}{FPGA} top-level files).
    \item \repopath{sdcard/}: SD card boot partition contents and root filesystem overlay, used to assemble the flashable image.
    \item \repopath{docs/}: pointer to the shared documentation folder (Chapter~\ref{chap:fpgalix-downloads}), where the reference manuals used throughout this work are kept.
    \item \repopath{WARNING.md}: camera hardware bug note (symlink to \repopath{hw/camera_adapter/WARNING.md}).
    \item \repopath{LICENSE}: MIT
\end{itemize}
